\documentclass[twoside,a5paper,11pt]{article} % nastaveni tak aby bylo pismo dostatecne velike
% ideální použití: 
% 1. vytvořit vlastní textový soubor, pojmenovat ho jménem zpěvníku "jmeno_zpevniku.txt" a vložit něj jednotlivé písně podle vzoru "\add{nazev_souboru_pisne.tex}"
%%% nazev_souboru je názem pod kterým je píseň ve správném formátu v kódování utf-8 uložená ve složce songs
% 2.  vložit soubor s písněmi sem pomocí "\input{nazev_meho_zpevniku.txt}", taky ve správném formátu a kódování utf-8
% 3. změnit následující příkaz tak, aby tam byla cesta ke složce, kde budou obrázky zpěvníku např: "\newcommand{\folder}{jmeno_zpevniku}"
\newcommand{\folder}{images}
% 4. do uvedené složky nahrát obrázky do zpěvníku - pouze některé písně mají v sobě obrázky, vkládají se jako "nazev_souboru_pisne.eps"


\usepackage[a5paper, top = 0.7cm, left = 1cm, right = 1cm, bottom = 1.5cm]{geometry} %velikost okrajů
\usepackage[utf8] {inputenc} %píšeme v utf-8, aby to pěkně jelo i linuxákům
\usepackage	[czech]	{babel}% nějaké další potřebné balíčky
\usepackage{guitar} %knihovna pro akordy
\usepackage{multicol} %knihovna pro více sloupců
\usepackage{graphicx}
\usepackage[T1]{fontenc}
\usepackage[unicode=true]{hyperref} %tvoří interaktivní obsah - kliknutím se dostaneš na danou písničku


\newcommand{\ch}[1]{\guitarChord{\textbf{#1}}}
\newcommand{\song}[1]{\clearpage\newpage \subsection{#1} \begin{enumerate}}    
\newcommand{\esong}{\end{enumerate}}
\newcommand{\ve}{\item }
\newcommand{\cho}[1]{\item[R:]#1}
\newcommand{\re}[1]{\item[#1]}
\newcommand{\add}[1]{\input{songs/#1}}
%\newcommand*\doublepage{  %vkládá se na začátek dlouhých písniček; zajistí, že bude začínat vždy na levé stránce, jinak doplní 1 bílou.
%  \clearpage
%  \ifodd\value{page}
%    \hbox{} \newpage
%  \fi
%}
\newcommand*\doublepage{\cleardoublepage} %v případě o 1 posunutých stránek

\widowpenalties 1 10000 %zakáže zlom stránky uprostřed sloky

\begin{document}

\author{Mara}%je vidět jen zde
\title{Zpěvník} %název se taky nikde nezobrazuje, ale pro pořádek




%\pagestyle{empty}
\tableofcontents %tady se generuje automatický obsah, tady se dá měnit. Správně se vygeneruje vždy až při druhém překladu.
\setcounter{secnumdepth}{1} % only chapter and sections will be numbered
\setcounter{tocdepth}{3} %== nastavení tak aby se zde v obsahu objevily pouze písně (ne sloky atd.)

\newpage
%\pagestyle{plain}     % zapne obyčejné číslování
%\setcounter{page}{6}  % nastaví čítač stránek znovu od jedné a můžeme začít vkládat písně
% překlad se musí nastavit na tex -> ps -> pdf

\input{songs/_seznam2.txt} % možnost přidání celého textového souboru s písněmi - více písní najednou
%\newpage
%\add{az_se_zepta_rano.tex} % jednotlivé písničky


%\newpage %začátek konce - tady začíná poslední strana
%\pagestyle{empty}
%\vspace*{\fill} %vertikální zarovnání na střed
%\begin{center} %horizontální zarovnání
%Pro tábor "šablona"\\
%v roce 2014\\
%kolektiv schopných autorů\\
%ve spolupráci s jazykem \TeX .\\
%Za skautské oddíly\\[2em] %vertikální mezera velikosti 2 písmen m
%3. roj světlušek "Světýlka"\\
%3. smečka vlčat\\
%20. oddíl "Králováci"\\
%\vspace*{\fill}

%\end{center}
\end{document}
